\Needspace{6\baselineskip}\Article{347}
\textsc{Theorem.} \textit{If \(F=\varphi(t,\allowbreak u.v\ldots)\) is a symmetric\HistoricalNote{*)}{\normalfont Functions are called symmetric, as is known, when all the indeterminates occur in them in the same way, or, more clearly, when they do not change however the indeterminates are permuted among themselves; such functions are \textit{for instance} the sum of all of them, the product of all of them, the sum of their pairwise products, etc.} algebraic} \textit{rational integral function of \(f\) indeterminates \(t,\allowbreak u,\allowbreak v\), etc., and, on substituting for these the \(f\) roots contained in the period \((f,\allowbreak \lambda)\), the value of \(F\) is reduced by the rules of Article 340 to the form} \[A+A^{\prime}[1]+A^{\prime\prime}[2]+\text{etc.}=W\] \textit{then the roots in this expression that belong to any one and the same period of \(f\) terms will have equal coefficients.}

\textit{Proof.} Let \([p],\allowbreak [q]\) be two roots belonging to one and the same period, and assume \(p,\allowbreak q\) positive and less than \(n\), so that we must show that \([p]\) and \([q]\) have the same coefficient in \(W\). Let \(q\equiv pg^{\nu e}\ (\mathrm{mod.}\ n)\); further, let the roots contained in \((f,\allowbreak \lambda)\) be \([\lambda],\allowbreak [\lambda^{\prime}],\allowbreak [\lambda^{\prime\prime}]\), etc., where we assume the numbers \(\lambda,\allowbreak \lambda^{\prime},\allowbreak \lambda^{\prime\prime}\), etc., positive and less than \(n\); finally, let the least positive residues of the numbers \(\lambda g^{\nu e},\allowbreak \lambda^{\prime} g^{\nu e},\allowbreak \lambda^{\prime\prime} g^{\nu e}\), etc., modulo \(n\), be \(\mu,\allowbreak \mu^{\prime},\allowbreak \mu^{\prime\prime}\), etc. These will clearly be identical with the numbers \(\lambda,\allowbreak \lambda^{\prime},\allowbreak \lambda^{\prime\prime}\), etc., although in a different order. Now Article 340 shows that \[\varphi([\lambda g^{\nu e}],[\lambda^{\prime} g^{\nu e}],[\lambda^{\prime\prime} g^{\nu e}]\ldots)=(\mathrm I)\] reduces to \[\begin{gathered}
A+A^{\prime}[g^{\nu e}]+A^{\prime\prime}[2g^{\nu e}]+\mathrm{etc.}\\
\text{or to}\quad A+A^{\prime}[\theta]+A^{\prime\prime}[\theta^{\prime}]+\mathrm{etc.}=(W^{\prime})
\end{gathered}\] where \(\theta,\allowbreak \theta^{\prime}\), etc., denote the least residues of the numbers \(g^{\nu e},\allowbreak 2g^{\nu e}\), etc., modulo \(n\). It is therefore clear that \([q]\) has the same coefficient in \((W^{\prime})\) as \([p]\) has in \((W)\). But it is readily seen that expanding the expression \((\mathrm I)\) gives the same result as expanding \(\varphi([\mu],\allowbreak [\mu^{\prime}].[\mu^{\prime\prime}]\ \text{etc.})\), since \(\mu\equiv\lambda g^{\nu e},\allowbreak \mu^{\prime}\equiv\lambda^{\prime} g^{\nu e}\), etc., \(\ (\mathrm{mod.}\ n)\); the latter expression in turn gives the same result as \(\varphi([\lambda],\allowbreak [\lambda^{\prime}],\allowbreak [\lambda^{\prime\prime}]\ \text{etc.})\), since the numbers \(\mu,\allowbreak \mu^{\prime},\allowbreak \mu^{\prime\prime}\), etc., differ from \(\lambda,\allowbreak \lambda^{\prime},\allowbreak \lambda^{\prime\prime}\), etc., only in their order, which makes no difference in a symmetric function. It follows that \(W^{\prime}\) will be completely identical with \(W\); hence the root \([q]\) will have the same coefficient in \(W\) as \([p]\). \textit{Q. E. D.}

It is therefore clear that \(W\) can be reduced to the form
\SourceLeaf{432}
\[A+a(f,1)+a^{\prime}(f,g)+a^{\prime\prime}(f,gg)\ldots+a^\varepsilon(f,g^{e-1})\] where the coefficients \(A,\allowbreak a\ldots a^\varepsilon\) are fixed quantities, which will also be integers if all the rational coefficients in \(F\) are integers. — Thus, \textit{for example}, if \(n=19\), \(f=6,\allowbreak \ \lambda=1\), and the function \(\varphi\) denotes the sum of the pairwise products of the indeterminates, its value reduces to \(3+(6,1)+(6,4)\).

It is further easily seen that, if we then substitute for \(t,\allowbreak u,\allowbreak v\), etc., the roots from another period \((f,k\lambda)\), the value of \(F\) becomes \[A+a(f,k)+a^{\prime}(f,kg)+a^{\prime\prime}(f,kgg)+\text{etc.}\]

\Needspace{6\baselineskip}\Article{348}
In any equation \[x^f-\alpha x^{f-1}+\beta x^{f-2}-\gamma x^{f-3}\ldots=0\] the coefficients \(\alpha,\allowbreak \beta,\allowbreak \gamma\), etc., are symmetric functions of the roots, namely \(\alpha\) the sum of all of them, \(\beta\) the sum of their pairwise products, \(\gamma\) the sum of their products taken three at a time, etc. Thus, in the equation whose roots are the roots contained in the period \((f,\lambda)\), the first coefficient will be \(=(f,\lambda)\), while each of the others can be reduced to a form such as \[A+a(f,1)+a^{\prime}(f,g)\ldots+a^\varepsilon(f,g^{e-1})\] where all of \(A,\allowbreak a,\allowbreak a^{\prime}\), etc., will be integers. It is also clear that the equation whose roots are those contained in any other period \((f,k\lambda)\) is derived from it by replacing, in each coefficient, \((f,1)\) by \((f,k)\); similarly, replacing \((f,g)\) by \((f,kg)\) and, in general, \((f,p)\) by \((f,kp)\). In this way we can therefore specify \(e\) equations \(z=0,\allowbreak z^{\prime}=0,\allowbreak z^{\prime\prime}=0\), etc., whose roots are the roots contained in \((f,1)\), in \((f,g)\), \((f,gg)\), etc., as soon as the \(e\) sums \((f,1),\allowbreak (f,g),\allowbreak (f,gg)\), etc., are known, or rather as soon as any \textit{one} of them has been found, since, by Article 346, all the others can be deduced rationally from one. At the same time, this resolves the function \(X\) into \(e\) factors of degree \(f\): for the product of the functions \(z,\allowbreak z^{\prime},\allowbreak z^{\prime\prime}\), etc., will clearly be \(=X\).

\textit{Example.} For \(n=19\), the sum of all the roots in the period \((6,1)\) is \(=(6,1)=\alpha\); the sum of their pairwise products is \(=3+(6,1)+(6,4)=\beta\); similarly,
\SourceLeaf{433}
the sum of the products taken three at a time is found to be \(=2+2(6,1)+(6,2)=\gamma\); the sum of the products taken four at a time \(=3+(6,1)+(6,4)=\delta\); the sum of the products taken five at a time \(=(6,1)=\varepsilon\); the product of all of them \(=1\). Thus the equation \[z=x^6-\alpha x^5+\beta x^4-\gamma x^3+\delta xx-\varepsilon x+1=0\] includes all the roots contained in \((6,1)\). If, in the coefficients \(\alpha,\allowbreak \beta,\allowbreak \gamma\), etc., we replace \((6,1),\allowbreak (6,2),\allowbreak (6,4)\) respectively by \((6,2),\allowbreak (6,4),\allowbreak (6,1)\), we obtain the equation \(z^{\prime}=0\), which includes the roots in \((6,2)\); if the same interchange is applied here again, we obtain the equation \(z^{\prime\prime}=0\), including the roots in \((6,4)\), and the product \(zz^{\prime}z^{\prime\prime}\) will be \(=X\).

\clearpage
\Needspace{6\baselineskip}\Article{349}
It is usually more convenient, especially when \(f\) is a large number, to deduce the coefficients \(\beta,\allowbreak \gamma\), etc., from the sums of the powers of the roots by Newton's theorem. For it is immediately clear that the sum of the squares of the roots contained in \((f,\lambda)\) is \(=(f,2\lambda)\), the sum of their cubes \(=(f,3\lambda)\), etc. Thus, writing \((f,\lambda),\allowbreak (f,2\lambda),\allowbreak (f,3\lambda)\), etc., as \(q,\allowbreak q^{\prime},\allowbreak q^{\prime\prime}\), etc., for brevity, we have \[\alpha=q,\quad 2\beta=\alpha q-q^{\prime},\quad 3\gamma=\beta q-\alpha q^{\prime}+q^{\prime\prime}\quad\text{etc.}\] Here the products of two periods must immediately be converted into sums of periods by Article 345. Thus, in our example, writing \((6,1),\allowbreak (6,2),\allowbreak (6,4)\) respectively as \(p,\allowbreak p^{\prime},\allowbreak p^{\prime\prime}\), we obtain \(q,\allowbreak q^{\prime},\allowbreak q^{\prime\prime},\allowbreak q^{\prime\prime\prime},\allowbreak q^{\prime\prime\prime\prime},\allowbreak q^{\prime\prime\prime\prime\prime}\) respectively \(=p,\allowbreak p^{\prime},\allowbreak p^{\prime},\allowbreak p^{\prime\prime},\allowbreak p^{\prime},\allowbreak p^{\prime\prime}\); hence \[\begin{aligned}
\alpha&=p,\quad2\beta=pp-p^{\prime}=6+2p+2p^{\prime\prime}\\
3\gamma&=(3+p+p^{\prime\prime})p-pp^{\prime}+p^{\prime}=6+6p+3p^{\prime}\\
4\delta&=(2+2p+p^{\prime})p-(3+p+p^{\prime\prime})p^{\prime}\\
&\qquad+pp^{\prime}-p^{\prime\prime}=12+4p+4p^{\prime\prime}\quad\mathrm{etc.}
\end{aligned}\] It suffices, however, to calculate only half the coefficients in this way. For it is not difficult to prove that the last ones, in reverse order, are either equal to the first ones, namely the last \(=1\), the next to last \(=\alpha\), the one before that \(=\beta\), etc., or can be deduced from them respectively by replacing \((f,1),\allowbreak (f,g)\), etc., by \((f,-1),\allowbreak (f,-g)\), etc., or \((f,n-1)\), \((f,n-g)\), etc. The first case occurs when \(f\) is even, the second when \(f\) is odd; the last coefficient, however, is always \(=1\). This rests on the theorem of Article 79; but, for brevity, we do not dwell on this subject.
\SourceLeaf{434}

\Needspace{6\baselineskip}\Article{350}
\textsc{Theorem.} \textit{Let \(n-1\) be the product of three positive integers \(\alpha,\allowbreak \beta,\allowbreak \gamma\); let} \textit{the period \((\beta\gamma,\lambda)\), which has \(\beta\gamma\) terms, consist of \(\beta\) smaller periods of \(\gamma\) terms,} \textit{namely \((\gamma,\lambda).(\gamma,\lambda^{\prime}),\allowbreak (\gamma,\lambda^{\prime\prime})\), etc. Suppose that, in a function of \(\beta\) indeterminates} \textit{of the same kind as in Article 347, namely \(F=\varphi(t,u,v\ldots)\), the indeterminates \(t,\allowbreak u,\allowbreak v\), etc.,} \textit{are replaced respectively by the sums \((\gamma,\lambda),\allowbreak (\gamma,\lambda^{\prime}),\allowbreak (\gamma,\lambda^{\prime\prime})\), etc., and its value is reduced by the rules of} \textit{Article 345, IV, to} \[\begin{aligned}
A+a(\gamma,1)+a^{\prime}(\gamma,g)&\ldots+a^\zeta(\gamma,g^{\alpha\beta-\alpha})\\
&\ldots+a^\theta(\gamma,g^{\alpha\beta-1})=W
\end{aligned}\] \textit{Then I say that, if \(F\) is a symmetric function, the periods in \(W\) that are contained in the same period} \textit{of \(\beta\gamma\) terms, \textit{that is}, generally such periods as \((\gamma,g^\mu)\) and \((\gamma,g^{\alpha\nu+\mu})\), where} \textit{\(\nu\) denotes any integer, will have the same coefficients.}

\textit{Proof.} Since the period \((\beta\gamma,\lambda g^\alpha)\) is identical with \((\beta\gamma,\lambda)\), the smaller periods \((\gamma,\lambda g^\alpha),\allowbreak (\gamma,\lambda^{\prime}g^\alpha),\allowbreak (\gamma,\lambda^{\prime\prime}g^\alpha)\), etc., of which the former clearly consists, must agree with those of which the latter consists, although in a different order. Thus, if we suppose that, on substituting these respectively for \(t,\allowbreak u,\allowbreak v\), etc., \(F\) becomes \(W^{\prime}\), \(W^{\prime}\) will coincide with \(W\). But by Article 347 we have \[\begin{aligned}
W^{\prime}&=A+a(\gamma,g^\alpha)+a^{\prime}(\gamma,g^{\alpha+1})\\
&\qquad\ldots+a^\zeta(\gamma,g^{\alpha\beta})\ldots+a^\theta(\gamma,g^{\alpha\beta+\alpha-1})\\
&=A+a(\gamma,g^\alpha)+a^{\prime}(\gamma,g^{\alpha+1})\\
&\qquad\ldots+a^\zeta(\gamma,1)\ldots+a^\theta(\gamma,g^{\alpha-1})
\end{aligned}\] Therefore, since this expression must agree with \(W\), the first, second, third, etc., coefficients in \(W\) (starting with \(a\)) must agree, respectively, with those in positions \(\alpha+1,\allowbreak \alpha+2,\allowbreak \alpha+3\), etc. It follows readily that, in general, the coefficients of the periods \((\gamma,g^\mu),\allowbreak (\gamma,g^{\alpha+\mu}),\allowbreak (\gamma,g^{2\alpha+\mu})\ldots(\gamma,g^{\nu\alpha+\mu})\), which occupy positions \(\mu+1,\allowbreak \alpha+\mu+1,\allowbreak 2\alpha+\mu+1\ldots\nu\alpha+\mu+1\), must agree with one another. \textit{Q. E. D.}

It is therefore clear that \(W\) can be reduced to the form \[A+a(\beta\gamma,1)+a^{\prime}(\beta\gamma,g)\ldots+a^\varepsilon(\beta\gamma,g^{\alpha-1})\] where all the coefficients \(A,\allowbreak a\), etc., will be integers if all the fixed coefficients in \(F\) are integers. It is further easily seen that, if we then substitute for the indeterminates in \(F\) the \(\beta\) periods of \(\gamma\) terms contained in another period of \(\beta\gamma\) terms, namely \((\beta\gamma,\lambda k)\), which are clearly \((\gamma,\lambda k),\allowbreak (\gamma,\lambda^{\prime}k),\allowbreak (\gamma,\lambda^{\prime\prime}k)\), etc., the resulting value will be \(A+a(\beta\gamma,k)+a^{\prime}(\beta\gamma,gk)\ldots+a^\varepsilon(\beta\gamma,g^{\alpha-1}k)\).
\SourceLeaf{435}

It is also clear that the theorem extends to the case where \(\alpha=1\), or \(\beta\gamma=n-1\): here \textit{all} the coefficients in \(W\) will be equal, so that \(W\) reduces to the form \(A+a(\beta\gamma,1)\).


